\documentclass{report}
\usepackage{amsmath,amssymb}
\setlength{\parindent}{0mm}
\setlength{\parskip}{1em}
\begin{document}
\begin{center}
\rule{6in}{1pt} \
{\large Jos� Mar�n \\
{\bf Preconditioners for nonsymmetric linear systems with low-rank skew-symmetric part}}

Institut de Matematica Multidisciplinar \\ Universitat Politecnica de Valencia \\ Cam� de Vera \\ 14 \\ 46022 Valencia \\ Spain
\\
{\tt jmarinma@imm.upv.es}\\
Juana Cerd�n\\
Danny Guerrero\\
Jos� Mas\end{center}

\newcommand{\reals}{\mathbb{R}}

In this work we study the iterative solution of nonsingular, nonsymmetric
linear systems of $n$ equations
\begin{equation} Ax=b \end{equation}
where the skew-symmetric part of the coefficient matrix $A$ can be
approximated by a low-rank matrix.
Consider $A=H+K$ where $H$ and $K$ are the symmetric and skew-symmetric
parts of $A$, respectively. It is assumed that the skew-symmetric matrix
can be written as $K=PQ^T + E$ for some full rank $P,Q \in \reals^{n
\times s}$ with $k \ll n$, and $\| E \| \ll 1$.

Different strategies have been proposed when the skew-symmetric part $K$
has exactly rank $s \ll n$. In \cite{Beckermann:2008} it is presented a
progressive GMRES method that allows for the short-term computation of an
orthogonal Krylov subspace basis. As pointed out in \cite{ESSSX.2012},
although the method is mathematically equivalent to full GMRES
\cite{sasc:86} in practice it may suffer from instabilities due to the
loss of orthogonality between the vectors of the generated Krylov
subspace basis. In the same paper, the authors propose a Schur complement
method that also permits the application of short-term formulas. The
method obtains an
approximate solution by applying the MINRES method $s+1$ times. The
authors also suggest that can be successfully applied as a preconditioner
for GMRES for the problem considered in this paper.

We study a method based on the framework proposed in \cite{cmm:16}.
Assuming that the matrix $H+PQ^T$ is nonsingular, our approach computes
an approximate LU factorization of the matrix
\begin{equation}\label{eq:augmented}
\left[
\begin{array}{cc}
H & P\\
-Q^T & I
\end{array}
\right]
\end{equation}
with the Balanced Incomplete Factorization (BIF) algorithm
\cite{bmmt:08,bmmt:10}. Interestingly, the matrix in (\ref{eq:augmented})
is
similar to the one used in \cite{ESSSX.2012} to develop the Schur
complement method, but in this work it is used to update a previously
computed preconditioner for the symmetric part $H$. Then, the
factorization is used as a preconditioner for the GMRES method. The
results of the numerical experiments for different problems will be
presented.

\begin{thebibliography}{1}

\bibitem{Beckermann:2008}
Bernhard Beckermann and Lothar Reichel.
\newblock The {A}rnoldi process and {GMRES} for nearly symmetric matrices.
\newblock {\em SIAM J. Matrix Anal. Appl.}, 30(1):102--120, February 2008.

\bibitem{bmmt:08}
R.~Bru, J.~Mar{\'{\i}}n, J.~Mas, and M.~T\accent23uma.
\newblock Balanced incomplete factorization.
\newblock {\em SIAM J. Sci. Comput.}, 30(5):2302--2318, 2008.

\bibitem{bmmt:10}
R.~Bru, J.~Mar{\'{\i}}n, J.~Mas, and M.~T\accent23uma.
\newblock Improved balanced incomplete factorization.
\newblock {\em SIAM J. Matrix Anal. Appl.}, 31(5):2431--2452, 2010.

\bibitem{cmm:16}
J.~Cerd\'an, J.~Mar\'{\i}n, and J.~Mas.
\newblock Low-rank updates of balanced incomplete factorization
preconditioners.
\newblock {\em Submitted.}

\bibitem{ESSSX.2012}
Mark Embree, Josef~A. Sifuentes, Kirk~M. Soodhalter, Daniel~B. Szyld, and Fei
Xue.
\newblock Short-term recurrence {K}rylov subspace methods for nearly hermitian
matrices.
\newblock {\em SIAM. J. Matrix Anal. and Appl.}, 33-2:480--500, 2012.

\bibitem{sasc:86}
Y.~Saad and M.~H. Schulz.
\newblock {GMRES}: {A} generalized minimal residual algorithm for solving
nonsymmetric linear systems.
\newblock {\em {SIAM} Journal on Scientific and Statistical Computing},
7:856--869, 1986.

\end{thebibliography}


\end{document}
